% Chapter: The global velocity.  Section 3 of the paper (sec:velocity).
% Owned shared labels: def:velocity:scale-family, def:velocity:placement,
% def:velocity:global-velocity.  Paper label: lem:global-velocity (Lemma 3.1 of the paper).

The local seed of the previous chapter is a single smooth, compactly supported,
divergence-free velocity $v $ on $\R^2 $ together with one axial wave number
$N=2\pi k_* $, one interval $I $ of admissible rescaled diffusivities, one contour
$\Gamma $ and the growth floor $g_*>0 $ (node \texttt{prop:local-seed}).  It produces
instability at one scale only.  This chapter installs rescaled copies of that seed at a
geometric sequence of horizontal scales $\ell_n\downarrow0 $, packs them into pairwise
disjoint discs of a fixed Euclidean chart of $\T^2 $, adds in each disc a small axial
velocity whose only role is to shift that cell's spectrum vertically by the fixed amount
$-i\tau $, and sums everything into one autonomous velocity field $u $ on $\T^3 $.

The point of the construction is a quantifier order: nothing here depends on the physical
diffusivity $\eps $.  In this blueprint's shorthand the order of the choices is
\[
  (v,I,\Gamma,\cG,g_*,f_*,L_*,q,\vartheta)
  \quad\longrightarrow\quad\tau
  \quad\longrightarrow\quad K_0 :
\]
the local seed, the single geometric ratio $q $ and the axial cut-off
$\vartheta $ come first, then the axial shift $\tau $ of the background chapter, then the
base Fourier index $K_0 $, which fixes the field $u $ and the threshold $\eps_0 $; only
afterwards is a cell matched to a given $\eps $.  The paper states this order in prose, in
Section~3.1 (``Choice of the parameters'') and in the paragraph collecting the parameter choices at
the end of Section~7.2, just before Remark~7.5 of the paper (\texttt{rem:flexible-buffer-scales});
it does not display the chain.  Every estimate that chooses $K_0 $ is uniform once the
earlier parameters are fixed.  The chapter ends
with the field's regularity (Lemma~\ref{lem:global-velocity}): $u $ is Lipschitz but not
$\C^1 $, and the failure of $\C^1 $ is exactly the accumulation of nonvanishing gradients at
infinitely many shrinking scales.  The dynamical consequences of the Lipschitz bound
(nodes \texttt{prop:particle-flow} and \texttt{cor:ideal-growth}) are in the particle
chapter.

Terminology: the paper calls the $y $-directions \emph{horizontal} and this chapter
follows it; the Lean identifiers use the adjective \texttt{transverse}
(\texttt{TransverseRealW1Inf}, \texttt{transverseTorusPoint},
\texttt{IsTransverseDivergenceFree}, \dots).  The correspondence is recorded once in
\texttt{FastDynamo/Conventions.lean}.

\section{Diffusivity intervals and horizontal scales}

\begin{definition}[The scale family]
  \label{def:velocity:scale-family}
  \lean{FastDynamo.ScaleRatios, FastDynamo.ScaleFamily,
    FastDynamo.ScaleFamily.fourierIndex, FastDynamo.ScaleFamily.ell,
    FastDynamo.ScaleFamily.bufferRadius, FastDynamo.ScaleFamily.cellCenter,
    FastDynamo.ScaleFamily.axialAmplitude,
    FastDynamo.ScaleFamily.diffusivityInterval, FastDynamo.ScaleFamily.epsilon0,
    FastDynamo.ScaleFamily.MatchingCell,
    FastDynamo.ScaleRatios.q_lt_scaleRatioCeiling, FastDynamo.ScaleFamily.rounding_bound,
    FastDynamo.ScaleFamily.TorusCellPlacement.FullVelocityConcrete.accumulationPoint}
  \leanok
  \uses{prop:local-seed, def:local:seed}
  Fix a ratio $1<q<\sqrt{\delta_+/\delta_-} $ and an integer $K_0 $ so large that
  $q+K_0^{-1}<\sqrt{\delta_+/\delta_-} $.  Set
  \[
    K_{n+1}=\lceil qK_n\rceil,
    \qquad
    \ell_n=\frac{k_*}{K_n},
    \qquad
    R_n=\sqrt{\ell_n},
    \qquad
    a_n=\frac{\tau}{2\pi K_n}=\frac{\tau}{2\pi k_*}\ell_n
  \]
  for a parameter $\tau\geq0 $ fixed later.  Here $\ell_n $ is the horizontal scale of the
  $n $-th rescaled copy of the local velocity (the \emph{core} scale), $R_n $ is the radius
  reserved around it (the \emph{buffer} scale) and $a_n $ is the amplitude of the added
  axial velocity.  The centres and the accumulation point are
  \[
    p_n=\Bigl(4\sum_{m\geq n}R_m,\tfrac12\Bigr),
    \qquad
    p_\infty=\bigl(0,\tfrac12\bigr),
  \]
  and the corresponding accumulation set in $\T^3 $ is the circle $\{p_\infty\}\times\T $.
  The diffusivity intervals and the threshold are
  \[
    \cI_n=\Bigl[\delta_-\frac{k_*^2}{K_n^2},\;\delta_+\frac{k_*^2}{K_n^2}\Bigr],
    \qquad
    \eps_0=\delta_+\frac{k_*^2}{K_0^2}.
  \]
  An index $n $ is \emph{admissible for} $\eps $ when $\delta_n:=\eps/\ell_n^2\in I $; an
  admissible index together with that membership is a \emph{matching cell}.
\end{definition}

In Lean the two rounds of choice are two structures: \texttt{ScaleRatios p} carries the
single ratio $q $ together with $q<\sqrt{\delta_+/\delta_-} $
(\texttt{ScaleRatios.q\_lt\_scaleRatioCeiling}), and \texttt{ScaleFamily r} carries only
$K_0 $ with the rounding bound $q+K_0^{-1}<\sqrt{\delta_+/\delta_-} $
(\texttt{ScaleFamily.rounding\_bound}), so that later stages may still enlarge $K_0 $:
the bound is monotone in $K_0 $, which is the paper's remark after
\texttt{eq:ratio-choice} that every subsequent enlargement of $K_0 $ preserves it
(\texttt{ScaleFamily.rounding\_bound\_of\_le}).
All of $K_n,\ell_n,R_n,a_n,p_n,\cI_n,\eps_0 $ are then functions of a
\texttt{ScaleFamily} (modules \texttt{FastDynamo/Global/ScaleSelection.lean} and
\texttt{FastDynamo/Global/CellGeometry.lean}); the point $p_\infty $ is the constant
\texttt{accumulationPoint} of
\texttt{FastDynamo/Global/TorusGlobalVelocityQualitative.lean}.  A
\texttt{MatchingCell f eps} is an index bundled with the proof that
$\eps/\ell_n^2\in I $; uniqueness of the index is neither claimed nor needed, exactly as
in the paper.

\textbf{Deviation: the free wave number $k_* $.}  The paper sets $\ell_n=1/K_n $, so that
$K_n\ell_n=1 $, $a_n=\tau\ell_n/(2\pi) $, $\cI_n=[\delta_-/K_n^2,\delta_+/K_n^2] $ and
$\eps_0=\delta_+/K_0^2 $.  Lean keeps the local axial wave number $N=2\pi k_* $ of the seed
as a parameter ($k_* $ a positive integer, the field \texttt{LocalSeedParameters.kStar}) and
sets $\ell_n=k_*/K_n $; the displays above are written in that generality.  The concrete
seed has $k_*=1 $ (\texttt{globalSeedInterface\_kStar}), where every formula of the definition is the
paper's.

\begin{lemma}[Exact identities and geometric decay of the scale family]
  \label{lem:velocity:scale-arithmetic}
  \lean{FastDynamo.ScaleFamily.adjacent_ratio_bounds,
    FastDynamo.ScaleFamily.adjacent_ratio_sq_lt,
    FastDynamo.ScaleFamily.axial_wavenumber_match,
    FastDynamo.ScaleFamily.bufferRadius_sq,
    FastDynamo.ScaleFamily.ell_div_bufferRadius,
    FastDynamo.ScaleFamily.axialAmplitude_match,
    FastDynamo.ScaleFamily.summable_bufferRadius,
    FastDynamo.ScaleFamily.tendsto_bufferRadius_div_ell_atTop,
    FastDynamo.ScaleFamily.epsilon0_formula,
    FastDynamo.ScaleFamily.iUnion_diffusivityInterval,
    FastDynamo.ScaleFamily.exists_matchingCell,
    FastDynamo.ScaleFamily.TorusCellPlacement.FullVelocityConcrete.tendsto_cellCenter_accumulationPoint}
  \leanok
  \uses{def:velocity:scale-family}
  For every scale family:
  \begin{enumerate}
    \item $q\leq K_{n+1}/K_n\leq q+1/K_0<\sqrt{\delta_+/\delta_-} $, the paper's
      \texttt{eq:ratio-choice};
    \item $2\pi K_n\ell_n=N $ and $2\pi K_na_n=\tau $;
    \item $R_n^2=\ell_n $ and $\ell_n/R_n=R_n $, hence $R_n/\ell_n\to\infty $;
    \item $\sum_{n\geq0}R_n<\infty $, and $p_n\to p_\infty $;
    \item $\bigcup_{n\geq0}\cI_n=(0,\eps_0] $ with $\eps_0=\delta_+k_*^2/K_0^2 $; hence
      every $0<\eps\leq\eps_0 $ admits at least one index $n $ with $\eps/\ell_n^2\in I $.
  \end{enumerate}
\end{lemma}

Items (i)--(v) are separate Lean theorems, each holding for every
\texttt{ScaleRatios} and \texttt{ScaleFamily}.  Item (i) is the full three-part chain
\texttt{eq:ratio-choice}, one Lean theorem,
\texttt{ScaleFamily.adjacent\_ratio\_bounds}.  In the paper its upper bound gives the
overlap of consecutive diffusivity intervals (Section~3.2), and its lower bound gives $r_m\geq q $
for $m<n $ and $r_m\leq q^{-1} $ for $m>n $ in the proof of Proposition~6.1 of the paper
(\texttt{prop:background-resolvent}), where the axial shifts of distinct cells are separated
through $c_q=(q-1)/(q+1) $.  Item (ii) is the exact frequency matching
that makes each cell a literal dilate of the local problem, together with the identity that
makes the added axial velocity act as the fixed spectral translation $-i\tau $ on the
fixed-$K_n $ sector.  The paper sets $\ell_n=1/K_n $ and uses $K_n\ell_n=1 $ directly,
without a label.  Lean keeps $k_* $ free, so \texttt{ScaleFamily.axial\_wavenumber\_match}
reads $2\pi K_n\ell_n=N=2\pi k_* $ and is the paper's identity at $k_*=1 $.

\begin{proof}
  \leanok
  \uses{def:velocity:scale-family}
  (i) is the rounding estimate $qK_n\leq\lceil qK_n\rceil<qK_n+1 $ combined with
  $K_0\leq K_n $ and the rounding bound $q+K_0^{-1}<\sqrt{\delta_+/\delta_-} $.  (ii) is
  immediate from $\ell_n=k_*/K_n $ and
  $a_n=\tau\ell_n/(2\pi k_*) $; (iii) from $R_n=\sqrt{\ell_n} $.  (iv) follows from the
  geometric bound $\ell_n\leq\ell_0q^{-n} $, hence $R_n\leq R_0q^{-n/2} $.  For (v), the
  ratio condition $(q+1/K_0)^2<\delta_+/\delta_- $
  (\texttt{ScaleFamily.adjacent\_ratio\_sq\_lt}) makes consecutive intervals overlap, and
  both endpoints tend to zero, so their union is the half-open interval down from $\eps_0 $.
\end{proof}

\begin{lemma}[The base index may be chosen above any later bound]
  \label{lem:velocity:base-index}
  \lean{FastDynamo.ScaleFamily.exists_ge, FastDynamo.exists_scale_selection,
    FastDynamo.ScaleFamily.exists_largeK0_chartPacking_and_supportFit,
    FastDynamo.ScaleFamily.rounding_bound_of_le,
    FastDynamo.ScaleFamily.TorusCellPlacement.nonempty}
  \leanok
  \uses{def:velocity:scale-family}
  Admissible ratios exist, and for every $M\in\Z_{\geq0} $ there is a scale family with
  $K_0\geq M $.  Moreover, for every $R_v\geq0 $ a large enough $K_0 $ satisfies both
  \[
    4\sum_{n\geq0}R_n+R_0<\tfrac14
    \qquad\text{and}\qquad
    R_v\ell_n<\frac{R_n}{10}\quad\text{for every }n .
  \]
\end{lemma}

The \texttt{exists\_ge} shape ($\forall M,\ \exists f,\ M\leq f.K_0 $) is what implements
the paper's dependency order: every later stage of the proof imposes one further
monotone lower bound on $K_0 $, and all of them are absorbed into a single maximum.  The
two displayed conditions are \texttt{eq:chart-packing} and
\texttt{eq:local-support-in-buffer}.  The single object that funnels every later bound
into one choice is \texttt{GlobalParameterChoice}, the node
\texttt{def:transfer:parameter-choice} of the transfer chapter.

\textbf{Deviation: explicit thresholds on $K_0 $.}  The paper asks for
\texttt{eq:chart-packing} ``for $K_0 $ sufficiently large'' and derives
\texttt{eq:local-support-in-buffer} from the explicit sufficient condition
$K_0>(10R_v)^2 $, via $R_v\ell_n/R_n=R_v/\sqrt{K_n}\leq R_v/\sqrt{K_0} $ (Section~3.3).  Lean
proves the same criterion in the form $10R_vR_0<1 $
(\texttt{ScaleFamily.supportScale\_lt\_tenth\_buffer}), which at $k_*=1 $ is exactly
$K_0>(10R_v)^2 $; the funnel
\texttt{ScaleFamily.exists\_ge\_chartPacking\_supportFit\_bufferRadius\_zero\_lt} behind
\texttt{exists\_largeK0\_chartPacking\_and\_supportFit} then chooses one explicit integer
$K_0>k_*/t^2 $ with
$t=\min\bigl(\tfrac14/(4/(1-q^{-1/2})+1),\ 1/(10(R_v+1)),\ \text{target}\bigr) $.  So the
packing threshold is the explicit $R_0\bigl(4/(1-q^{-1/2})+1\bigr)<\tfrac14 $ (via
\texttt{chartPacking\_of\_geometric\_bound}), and the support-fit threshold used
is $K_0>k_*(10(R_v+1))^2 $, slightly larger than the paper's $(10R_v)^2 $ (the shift by one
keeps the bound meaningful at $R_v=0 $).  Both are sufficient conditions of the paper's kind;
nothing downstream depends on the particular values.

\begin{proof}
  \leanok
  \uses{def:velocity:scale-family, lem:velocity:scale-arithmetic}
  Both conditions are controlled by $R_0=\sqrt{k_*/K_0} $ alone.  By the geometric bound
  $R_n\leq R_0(\sqrt{q^{-1}})^n $ the packing quantity is at most
  $R_0\bigl(4/(1-\sqrt{q^{-1}})+1\bigr) $, and $R_v\ell_n<R_n/10 $ reduces to
  $10R_vR_0<1 $ because $\ell_n/R_n=R_n\leq R_0 $.  Choosing $K_0 $ beyond the corresponding
  threshold, and beyond the prescribed $M $, gives both.
\end{proof}

\section{Disjoint placement and added axial velocity profiles on $\T^2$}

\begin{definition}[Placement of the cells in the chart and in the torus]
  \label{def:velocity:placement}
  \lean{FastDynamo.ScaleFamily.placementChart, FastDynamo.ScaleFamily.bufferDisk,
    FastDynamo.ScaleFamily.localCoreDisk,
    FastDynamo.ScaleFamily.TorusCellPlacement,
    FastDynamo.transverseTorusPoint,
    FastDynamo.ScaleFamily.TorusCellPlacement.torusBufferDisk,
    FastDynamo.ScaleFamily.TorusCellPlacement.torusLocalCoreDisk,
    FastDynamo.GlobalSeedInterface.placementRadius,
    FastDynamo.GlobalSeedInterface.profileSupportRadius_le_placementRadius,
    FastDynamo.GlobalSeedInterface.witnessSupportRadius_le_placementRadius,
    FastDynamo.ScaleFamily.TorusCellPlacement.restrictSupportRadius}
  \leanok
  \uses{def:velocity:scale-family, lem:velocity:base-index}
  All cells are placed in one fixed Euclidean coordinate square of $\T^2 $.  The $n $-th
  \emph{buffer disc} is $D(p_n,R_n) $ --- what the paper calls the $n $-th \emph{cell}
  --- and the $n $-th \emph{local-flow core} --- the paper's \emph{core} --- is the disc
  $D(p_n,R_v\ell_n) $ carrying the support of the rescaled reference profile, where $R_v $ is
  chosen with
  \[
    R_v\geq L_*
    \qquad\text{and}\qquad
    \supp v\cup\supp\nabla_\perp v\subset D(0,R_v)
  \]
  (Section~3.3 of the paper).  The corresponding subsets
  of $\T^3 $ are $D(p_n,R_n)\times\T $ and $D(p_n,R_v\ell_n)\times\T $.  A \emph{placement}
  is a scale family together with the two certificates of
  Lemma~\ref{lem:velocity:base-index}: chart packing and support fit.
\end{definition}

In Lean the chart is \texttt{ScaleFamily.placementChart}, the discs are
\texttt{bufferDisk} and \texttt{localCoreDisk} (Euclidean, in $\R\times\R $), and the
placement certificate is the structure \texttt{ScaleFamily.TorusCellPlacement
supportRadius}, whose two fields are exactly \texttt{chartPacking} and
\texttt{supportFit} (module \texttt{FastDynamo/Global/TorusCellPlacement.lean}).  The
quotient map $\R^2\to\T^2 $ is \texttt{transverseTorusPoint}
(\texttt{FastDynamo/Analysis/TransverseChart.lean}), and \texttt{torusBufferDisk},
\texttt{torusLocalCoreDisk} are the images of the Euclidean discs
(\texttt{FastDynamo/Global/TorusChartEmbedding.lean}).

Two differences from the paper.  First, the Lean chart is
$(0,\tfrac14)\times(\tfrac14,\tfrac34) $, the right half of the paper's
$(-\tfrac14,\tfrac14)\times(\tfrac14,\tfrac34) $; since every centre has positive first
coordinate and $p_\infty $ is the left edge, every disc still lies in it and nothing about
the accumulation changes.  Second, the concrete local profile is the smoothed Ponomarenko
velocity of the seed chapter, whose own support radius is $1+h $ for the smoothing width
$h $; nothing in the development bounds the witness radius $L_* $ above, so the paper's
$R_v $ is realised as
\[
  R_v=\texttt{GlobalSeedInterface.placementRadius}=\max(1+h,L_*),
\]
which is the smallest choice satisfying both requirements of the definition.  The two
inequalities are the Lean theorems
\texttt{GlobalSeedInterface.profileSupportRadius\_le\_placementRadius} and
\texttt{GlobalSeedInterface.witnessSupportRadius\_le\_placementRadius}, the second being
exactly $R_v\geq L_* $.  The placement is carried at this joint radius (the record field
\texttt{GlobalParameterChoice.jointPlacement}) and read back at $1+h $ for the velocity
constructors through \texttt{TorusCellPlacement.restrictSupportRadius}, which leaves the
scale family untouched; that reading is the derived accessor
\texttt{GlobalParameterChoice.placement}.  The cost of the joint radius is nothing: it is
one call of the funnel of Lemma~\ref{lem:velocity:base-index} instead of two, and it is what
makes \texttt{eq:scaled-witness-support} a consequence rather than a further threshold
(node \texttt{lem:transfer:witness}).

\begin{lemma}[The placed cells are disjoint, in the chart and in the torus]
  \label{lem:velocity:disjoint-placement}
  \lean{FastDynamo.ScaleFamily.bufferDisk_subset_placementChart,
    FastDynamo.ScaleFamily.disjoint_bufferDisk,
    FastDynamo.ScaleFamily.localCoreDisk_subset_bufferDisk,
    FastDynamo.transverseTorusPoint_injectiveOn_placementChart,
    FastDynamo.ScaleFamily.TorusCellPlacement.disjoint_torusBufferDisk,
    FastDynamo.ScaleFamily.TorusCellPlacement.torusLocalCoreDisk_subset_torusBufferDisk,
    FastDynamo.ScaleFamily.TorusCellPlacement.exists_matchingCell_with_torusPlacedCore}
  \leanok
  \uses{def:velocity:placement}
  For a placement: every $D(p_n,R_n) $ lies in the fixed chart; the discs $D(p_n,R_n) $ are
  pairwise disjoint; and $D(p_n,R_v\ell_n)\subset D(p_n,R_n) $ for every $n $.  The quotient
  map is injective on the chart, so the images in $\T^2 $ are again pairwise disjoint, with
  each core image inside its own buffer image.  Consequently at most one cell is active at
  each point of $\T^2 $.
\end{lemma}

\begin{proof}
  \leanok
  \uses{def:velocity:placement, lem:velocity:base-index}
  Chart containment: the first coordinate of a point of $D(p_n,R_n) $ lies between
  $4\sum_{m\geq n}R_m-R_n>0 $ and $4\sum_{m\geq0}R_m+R_0<\tfrac14 $, while its second
  coordinate differs from $\tfrac12 $ by less than $R_n\leq R_0<\tfrac14 $.  Disjointness:
  the centres are collinear with $\abs{p_n-p_{n+1}}=4R_n>R_n+R_{n+1} $, and summing adjacent
  gaps gives $\abs{p_n-p_m}>R_n+R_m $ for $n<m $.  The core inclusion is the support fit
  $R_v\ell_n<R_n/10 $.  Injectivity on the chart is uniqueness of representatives in
  $[0,1) $ for $\R/\Z $, applied coordinatewise; it transports both statements to the torus.
\end{proof}

\begin{lemma}[The added axial velocity and its cut-off gradient]
  \label{lem:velocity:axial-plateau}
  \lean{FastDynamo.ScaleFamily.TorusCellPlacement.axialPlateauBump,
    FastDynamo.ScaleFamily.TorusCellPlacement.axialPlateau_eq_one_of_norm_pairToWholePlane_le,
    FastDynamo.ScaleFamily.TorusCellPlacement.axialOffsetEuclidean,
    FastDynamo.ScaleFamily.TorusCellPlacement.axialOffsetEuclidean_eq_axialAmplitude_of_core,
    FastDynamo.ScaleFamily.TorusCellPlacement.axialOffsetEuclidean_eq_zero_of_not_mem_bufferDisk,
    FastDynamo.ScaleFamily.TorusCellPlacement.norm_axialOffsetEuclideanDerivative_le,
    FastDynamo.ScaleFamily.axialGradientScale,
    FastDynamo.ScaleFamily.axialGradientScale_formula,
    FastDynamo.ScaleFamily.axialGradientScale_le_zero}
  \leanok
  \uses{def:velocity:placement, lem:velocity:scale-arithmetic}
  Let $\vartheta\in\C^\infty_c([0,\infty)) $ equal one on $[0,3/4] $ with
  $\supp\vartheta\subset[0,1) $, and put
  \[
    u_n^\mathrm{ax}(y)=a_n\vartheta\Bigl(\frac{\abs{y-p_n}}{R_n}\Bigr)e_3 .
  \]
  Then $u_n^\mathrm{ax}=a_ne_3 $ on the \emph{constant axial region} $D(p_n,3R_n/4) $,
  $\supp u_n^\mathrm{ax}\subset D(p_n,R_n) $, the field is independent of $x_3 $ and
  divergence-free, and
  \[
    \norm{\nabla u_n^\mathrm{ax}}_\infty\leq\norm{\vartheta'}_\infty\frac{a_n}{R_n},
    \qquad
    \frac{a_n}{R_n}=\frac{\tau}{N}R_n=\frac{\tau}{2\pi\sqrt{k_*K_n}}
    \leq\frac{a_0}{R_0} .
  \]
  In particular the supremum
  $\eta_\mathrm{ax}=\sup_n\norm{\vartheta'}_\infty a_n/R_n $ is attained at $n=0 $ and tends to
  $0 $ as $K_0\to\infty $.
\end{lemma}

The Lean plateau is a \texttt{ContDiffBump} at the origin of the pair chart, equal to one
on the closed ball of radius $3/4 $ and vanishing outside the open unit ball;
$\norm{\vartheta'}_\infty $ is realised by its Lipschitz constant
\texttt{axialPlateauLipschitzConstant}, and the quantity $a_n/R_n $ is
\texttt{ScaleFamily.axialGradientScale}.  The closed expression
$\eta_\mathrm{ax}=\norm{\vartheta'}_\infty\tau/(2\pi\sqrt{k_*K_0}) $ is not a single Lean
lemma: what is proved is the monotonicity $a_n/R_n\leq a_0/R_0 $ together with
$a_0/R_0=(\tau/N)R_0 $ and $R_0=\sqrt{k_*/K_0} $, which is the same statement after
substitution.  The paper uses $\eta_\mathrm{ax} $ only as a gradient bound (in the proof of
Lemma~3.1 of the paper and in \texttt{eq:axial-shear-bound} of the proof of
Proposition~4.1 of the paper); its smallness for large $K_0 $ is noted in Section~3.3 but is not one
of the thresholds on $K_0 $, neither in the paper nor in
\texttt{GlobalParameterChoice}.  The axial taper is \emph{radial} --
$\vartheta(\abs{y-p_n}/R_n)$ in the Euclidean metric of the chart, not in the sup metric --
which is what makes the plateau $D(p_n,3R_n/4)$ a Euclidean disc and lets the support fit of
the placement be read off without losing a factor two.

\begin{proof}
  \leanok
  \uses{def:velocity:placement, lem:velocity:scale-arithmetic}
  The plateau statements are the defining properties of $\vartheta $, transported by
  $y\mapsto(y-p_n)/R_n $.  The chain rule produces the factor $R_n^{-1} $ in the gradient,
  and $a_n/R_n=(\tau/N)\ell_n/R_n=(\tau/N)R_n $ by
  Lemma~\ref{lem:velocity:scale-arithmetic}(iii); monotonicity in $n $ follows from
  $R_n\leq R_0 $.  Divergence-freeness holds because the only nonzero component is the third
  one and it does not depend on $x_3 $.
\end{proof}

\begin{definition}[The global velocity]
  \label{def:velocity:global-velocity}
  \lean{FastDynamo.ScaleFamily.TorusCellPlacement.parameterizedHelicalCellSource,
    FastDynamo.ScaleFamily.TorusCellPlacement.FullVelocityConcrete.parameterizedFullEuclideanCellRealSource,
    FastDynamo.ScaleFamily.TorusCellPlacement.FullVelocityConcrete.parameterizedFullGlobalEuclideanRealSource,
    FastDynamo.ScaleFamily.TorusCellPlacement.FullVelocityConcrete.parameterizedFullGlobalPeriodizedRealSource,
    FastDynamo.ScaleFamily.TorusCellPlacement.FullVelocityConcrete.parameterizedFullGlobalRealTorusSource,
    FastDynamo.ScaleFamily.TorusCellPlacement.FullVelocityConcrete.parameterizedFullGlobalW1Inf,
    FastDynamo.TransverseRealW1Inf,
    FastDynamo.literalVelocity}
  \leanok
  \uses{def:velocity:placement, lem:velocity:disjoint-placement, lem:velocity:axial-plateau,
    def:local:seed}
  With $u_n^\mathrm{loc}(y)=\ell_nv\bigl((y-p_n)/\ell_n\bigr) $ and $u_n^\mathrm{ax} $ as above,
  the \emph{global velocity} is
  \[
    u(y,x_3)=\sum_{n\geq0}\bigl(u_n^\mathrm{loc}(y)+u_n^\mathrm{ax}(y)\bigr).
  \]
  It is independent of $x_3 $, and by Lemma~\ref{lem:velocity:disjoint-placement} at most
  one summand is nonzero at each point, so the sum is defined pointwise without any
  convergence hypothesis.
\end{definition}

Lean builds this in three stages (modules
\texttt{FastDynamo/Global/TorusVelocityConcrete.lean},
\texttt{FastDynamo/Global/TorusAxialVelocityConcrete.lean},
\texttt{FastDynamo/Global/TorusFullVelocityConcrete.lean},
\texttt{FastDynamo/Global/TorusGlobalVelocityLipschitz.lean}):
\begin{itemize}
  \item the \emph{chart-side representative}
    \texttt{parameterizedFullGlobalEuclideanRealSource}, the literal pointwise sum
    $\sum_n(\re u_n^\mathrm{loc}+u_n^\mathrm{ax}) $ on $\R^2 $, supported in the placement
    chart;
  \item its \emph{$\Z^2 $-periodisation}
    \texttt{parameterizedFullGlobalPeriodizedRealSource}, which again has at most one
    nonzero translate at each point of the plane and descends to the torus-side
    representative \texttt{parameterizedFullGlobalRealTorusSource};
  \item the \emph{torus-side $\W^{1,\infty} $ class} \texttt{parameterizedFullGlobalW1Inf},
    an element of \texttt{TransverseRealW1Inf}: a real $\Lp^\infty $ field on $\T^2 $
    bundled with two $\Lp^\infty $ derivative fields subject to the Fourier
    weak-derivative identities
    $\widehat{\partial_ju}(m)=2\pi i\,m_j\,\widehat u(m) $.  Membership in
    $\W^{1,\infty} $ is thus carried by the type; the class is the sum of a pointwise
    disjoint-support limit (the helical part, whose cell derivatives are \emph{not}
    $\Lp^\infty $-summable) and an absolutely convergent $\Lp^\infty $ series
    (the axial part).
\end{itemize}
The axial shift is the parameter \texttt{tau}; the value used downstream is the chain's
\texttt{axialShift}.  The velocity of the main theorem is
\texttt{GlobalParameterChoice.velocity}, that is
\texttt{parameterizedFullGlobalW1Inf} with the helical
speeds pinned to the seed's own $(\Omega,U) $, so that the on-core identity of
Lemma~\ref{lem:velocity:matching-cell} produces the seed velocity $v $.  Finally
\texttt{literalVelocity} is the axially independent extension of the torus-side
representative to $\T^3 $, which is the field the paper writes as $u $.

\begin{lemma}[On a cell core the global field is a rescaled seed copy plus $a_ne_3 $]
  \label{lem:velocity:matching-cell}
  \lean{FastDynamo.ScaleFamily.TorusCellPlacement.FullVelocityConcrete.parameterizedFullGlobalTorusSource_eq_rescaledLocalCopy_add_axialShift_of_mem_core,
    FastDynamo.ScaleFamily.TorusCellPlacement.FullVelocityConcrete.parameterizedFullGlobalTorusFirstDerivative_eq_rescaledLocalCopy_of_mem_core,
    FastDynamo.ScaleFamily.TorusCellPlacement.FullVelocityConcrete.parameterizedFullGlobalRealTorusSource_eq_axialShift_of_mem_plateau_of_not_mem_core}
  \leanok
  \uses{def:velocity:global-velocity, lem:velocity:disjoint-placement,
    lem:velocity:axial-plateau, lem:velocity:scale-arithmetic}
  For every $n $ and every $y $ in the $n $-th local-flow core,
  \[
    u(y)=\ell_nv\Bigl(\frac{y-p_n}{\ell_n}\Bigr)+a_ne_3,
    \qquad
    \nabla_\perp u(y)=(\nabla_\perp v)\Bigl(\frac{y-p_n}{\ell_n}\Bigr),
  \]
  and for $y\in D(p_n,3R_n/4) $ outside the core, $u(y)=a_ne_3 $.  Since $2\pi K_na_n=\tau $,
  the added axial velocity acts on the fixed-$K_n $ Fourier subspace
  $B(y,x_3)=\e^{2\pi iK_nx_3}b(y) $ as the constant translation
  $-(a_ne_3\cdot\nabla)B=-i\tau B $, which changes no spectral real part.
\end{lemma}

The last two Lean names state the same identities almost everywhere for the stored
$\Lp^\infty $ classes at the selected matching cell, in the form the parametrix chapter
consumes (module \texttt{FastDynamo/Global/TorusSelectedCellChartAgreement.lean}); there
the local profile is exactly the seed's, because the helical speeds of the global
velocity are pinned to the seed's profile data (the fields \texttt{smoothingWidth},
\texttt{angularSpeed} and \texttt{axialSpeed} of \texttt{GlobalSeedInterface}).  The consequent operator identity
$U_n^{-1}A_nU_n=L_{\delta_n}-i\tau $ is the exact-renormalisation node
\texttt{lem:exact-renormalization} of the operator chapter (Lemma~5.2 of the paper,
``Scaling'', identity \texttt{eq:scaling}), which is where the paper carries out the
$-i\tau $ computation; it is not restated here.

\begin{proof}
  \leanok
  \uses{lem:velocity:disjoint-placement, lem:velocity:axial-plateau,
    lem:velocity:scale-arithmetic}
  By disjointness only the $n $-th summand is active on $D(p_n,R_n) $.  The support fit puts
  the core well inside $D(p_n,3R_n/4) $, where the axial cut-off is identically one, so the
  axial summand is the constant $a_ne_3 $ there and contributes nothing to the derivative.
  The helical summand is by definition $\ell_nv((y-p_n)/\ell_n) $, whose derivative is
  $(\nabla_\perp v)((y-p_n)/\ell_n) $ because the amplitude $\ell_n $ exactly cancels the
  dilation.  Outside the core the helical summand vanishes.  The spectral statement is the
  identity $2\pi K_na_n=\tau $ of Lemma~\ref{lem:velocity:scale-arithmetic}(ii).
\end{proof}

\section{Regularity and incompressibility of the velocity field on $\T^3$}

\begin{lemma}[Reality, $x_3 $-independence and incompressibility]
  \label{lem:velocity:divergence-free}
  \lean{FastDynamo.IsTransverseDivergenceFree,
    FastDynamo.ScaleFamily.TorusCellPlacement.FullVelocityConcrete.parameterizedFullGlobalTorusSource_pointwiseReal,
    FastDynamo.ScaleFamily.TorusCellPlacement.FullVelocityConcrete.parameterizedFullGlobalTorusFirstDerivative_transverseDivergence,
    FastDynamo.ScaleFamily.TorusCellPlacement.FullVelocityConcrete.parameterizedFullGlobalW1Inf_isTransverseDivergenceFree}
  \leanok
  \uses{def:velocity:global-velocity}
  The global velocity is real valued, independent of $x_3 $, and $\Div u=0 $: pointwise
  $\partial_1u_1+\partial_2u_2=0 $ for the displayed representatives, hence $\Div u=0 $ in
  the $\W^{1,\infty} $ graph sense.
\end{lemma}

In Lean, $x_3 $-independence is by type (the class lives on the transverse torus $\T^2 $),
reality is the vanishing of the imaginary part of every component of the complexified
representative, and $\Div u=0 $ is the predicate \texttt{IsTransverseDivergenceFree}: the
vanishing in $\Lp^\infty $ of the sum of the two stored diagonal derivative components.
Because the stored derivatives satisfy the Fourier weak-derivative identities, this is the
paper's distributional statement and not merely a pointwise one.

\begin{proof}
  \leanok
  \uses{def:velocity:global-velocity, lem:velocity:disjoint-placement,
    lem:velocity:axial-plateau}
  Each rescaled local copy is divergence-free because $\Div_\perp v_\perp=0 $ and dilation
  preserves that identity; each added axial velocity is divergence-free because it is
  parallel to $e_3 $ and independent of $x_3 $.  The supports are pairwise disjoint, so the
  identity for the sum is the identity for the unique active summand at each point.  Passing
  from the pointwise identity to the graph identity uses that the stored derivative fields
  are represented almost everywhere by the displayed pointwise sums.
\end{proof}

\begin{lemma}[Explicit uniform derivative bound and Lipschitz constant]
  \label{lem:velocity:w1inf-bound}
  \lean{FastDynamo.ScaleFamily.TorusCellPlacement.FullVelocityConcrete.norm_parameterizedFullEuclideanCellRealSource_le,
    FastDynamo.ScaleFamily.TorusCellPlacement.FullVelocityConcrete.globalVelocityCoordinateDerivativeBound,
    FastDynamo.ScaleFamily.TorusCellPlacement.FullVelocityConcrete.globalVelocityLipschitzBound,
    FastDynamo.ScaleFamily.TorusCellPlacement.FullVelocityConcrete.norm_fderiv_parameterizedFullGlobalEuclideanRealSource_le,
    FastDynamo.ScaleFamily.TorusCellPlacement.FullVelocityConcrete.lipschitzWith_parameterizedFullGlobalEuclideanRealSource,
    FastDynamo.ScaleFamily.TorusCellPlacement.FullVelocityConcrete.lipschitzWith_parameterizedFullGlobalRealTorusSource,
    FastDynamo.ScaleFamily.TorusCellPlacement.FullVelocityConcrete.lipschitzWith_literalVelocity_parameterizedFullGlobalRealTorusSource}
  \leanok
  \uses{def:velocity:global-velocity, lem:velocity:axial-plateau,
    lem:velocity:disjoint-placement}
  Each combined cell obeys
  $\norm{u_n^\mathrm{loc}+u_n^\mathrm{ax}}_\infty\leq(\norm v_\infty+\tau/N)\ell_n $, and there is
  an explicit constant
  \[
    L=2\Bigl(\norm{\nabla v}_\infty+\norm{\vartheta'}_\infty\frac{a_0}{R_0}\Bigr)
  \]
  with $\norm{Du(x)}\leq L $ at every point of the chart.  Consequently $u $ is Lipschitz
  with constant $2L $ on $\T^2 $ and, after the axial extension, on $\T^3 $.
\end{lemma}

This is the quantitative half of the paper's $\W^{1,\infty} $ claim, in slightly
different bookkeeping.  The proof of Lemma~3.1 of the paper bounds
$\norm{Du}\leq\norm{\nabla v}_\infty+\eta_\mathrm{ax} $ at every point; Lean proves
an everywhere bound on the chart carrying two extra factors of $2 $: one because the
operator norm on the pair chart is estimated by the sum of the two coordinate derivatives,
one because the periodisation estimate compares two competing integer translates.  The
role of $\norm{\nabla v}_\infty $ is played by the sum of the two first-order Schwartz
seminorms of the smoothed profile.  The constant is explicit but not claimed optimal, and
no later step uses its value.

\begin{proof}
  \leanok
  \uses{def:velocity:global-velocity, lem:velocity:axial-plateau,
    lem:velocity:disjoint-placement}
  Directly from the definitions, $\norm{u_n^\mathrm{loc}}_\infty\leq\ell_n\norm v_\infty $ and
  $\norm{u_n^\mathrm{ax}}_\infty\leq a_n=(\tau/N)\ell_n $.  For the derivative, the amplitude
  $\ell_n $ of the local copy exactly cancels its dilation, so
  $\norm{\nabla u_n^\mathrm{loc}}_\infty\leq\norm{\nabla v}_\infty $ uniformly in $n $, while
  $\norm{\nabla u_n^\mathrm{ax}}_\infty\leq\norm{\vartheta'}_\infty a_n/R_n
  \leq\norm{\vartheta'}_\infty a_0/R_0 $ by Lemma~\ref{lem:velocity:axial-plateau}.  Since
  the supports are pairwise disjoint, only one cell contributes at each point, which gives
  the bound for the whole sum; the mean value inequality turns it into a Lipschitz
  statement, and the descent lemma for $\Z^2 $-periodic Lipschitz maps transports it to
  $\T^2 $ and then to $\T^3 $.
\end{proof}

\begin{lemma}[Quadratic decay and differentiability at the accumulation circle]
  \label{lem:velocity:accumulation-decay}
  \lean{FastDynamo.ScaleFamily.TorusCellPlacement.FullVelocityConcrete.three_mul_bufferRadius_lt_dist_accumulationPoint_of_mem_bufferDisk,
    FastDynamo.ScaleFamily.TorusCellPlacement.FullVelocityConcrete.parameterizedFullGlobalEuclideanRealSource_eq_cell_of_mem_bufferDisk,
    FastDynamo.ScaleFamily.TorusCellPlacement.FullVelocityConcrete.parameterizedFullGlobalEuclideanRealSource_eq_zero_of_outside_buffers,
    FastDynamo.ScaleFamily.TorusCellPlacement.FullVelocityConcrete.norm_parameterizedFullGlobalEuclideanRealSource_le_dist_sq,
    FastDynamo.ScaleFamily.TorusCellPlacement.FullVelocityConcrete.parameterizedFullGlobalEuclideanRealSource_hasFDerivAt_accumulationPoint,
    FastDynamo.ScaleFamily.TorusCellPlacement.FullVelocityConcrete.parameterizedFullGlobalEuclideanRealSource_differentiableAt_accumulationPoint}
  \leanok
  \uses{def:velocity:global-velocity, lem:velocity:disjoint-placement,
    lem:velocity:scale-arithmetic}
  With $C=\norm v_\infty+\tau/N $,
  \[
    \abs{u(y,x_3)}\leq C\abs{y-p_\infty}^2 ,
  \]
  so $u $ vanishes at $p_\infty $ and is differentiable there with derivative zero.  The
  bound does not involve $x_3 $, so the same holds at every point of
  $\{p_\infty\}\times\T $.
\end{lemma}

The Lean statements are for the chart-side representative.  The only geometric inputs are
$\ell_n=R_n^2 $ and the separation $\abs{y-p_\infty}\geq3R_n $ for $y\in D(p_n,R_n) $.

\begin{proof}
  \leanok
  \uses{def:velocity:global-velocity, lem:velocity:disjoint-placement,
    lem:velocity:scale-arithmetic}
  If $y\in D(p_n,R_n) $ then
  $\abs{y-p_\infty}\geq\abs{p_n-p_\infty}-R_n=4\sum_{m\geq n}R_m-R_n\geq3R_n $, and there the
  field is the single $n $-th cell, of size at most $C\ell_n=CR_n^2 $; hence
  $\abs{u(y)}\leq C\abs{y-p_\infty}^2/9\leq C\abs{y-p_\infty}^2 $.  Outside all buffer discs
  the field vanishes, so the bound holds everywhere.  A field that is
  $O(\abs{y-p_\infty}^2) $ at $p_\infty $ is Fr\'echet differentiable there with derivative
  zero.
\end{proof}

\begin{lemma}[The global velocity is non-zero for a positive axial shift]
  \label{lem:velocity:nonvanishing}
  \lean{FastDynamo.ScaleFamily.TorusCellPlacement.FullVelocityConcrete.parameterizedFullGlobalTorusSource_ne_zero,
    FastDynamo.ScaleFamily.TorusCellPlacement.FullVelocityConcrete.parameterizedFullGlobalW1Inf_ne_zero}
  \leanok
  \uses{def:velocity:global-velocity, lem:velocity:matching-cell, lem:velocity:axial-plateau}
  If $\tau>0 $ then $u\neq0 $: on the nonempty set
  $D(p_n,3R_n/4)\setminus\supp u_n^\mathrm{loc} $ the field equals $a_ne_3\neq0 $.  The axial
  shift used in the construction is strictly positive, so the velocity of the main
  theorem is non-zero.
\end{lemma}

Because the torus-side object is an $\Lp^\infty $ equivalence class, the Lean proof needs
a set of positive measure rather than one point: it uses the explicit annular witness at
distance $R_n/2 $ from $p_n $ (outside the core by the support fit, inside the constant
axial plateau) together with a ball of radius $R_n/20 $ around it.  At $\tau=0 $ nothing
is proved; a helical-core witness would be needed, and the library contains none.
This is harmless: the shift used downstream is strictly positive.

\begin{proof}
  \leanok
  \uses{lem:velocity:matching-cell, lem:velocity:axial-plateau}
  Take the point at distance $R_n/2 $ from $p_n $.  The support fit $R_v\ell_n<R_n/10 $ puts
  it outside the local-flow core, and $R_n/2<3R_n/4 $ puts it inside the constant axial
  region, so by Lemma~\ref{lem:velocity:matching-cell} the field equals $a_ne_3 $ there,
  with $a_n=\tau\ell_n/(2\pi k_*)>0 $.  A ball of radius $R_n/20 $ around that point stays
  in the same region and has positive measure, which upgrades the pointwise statement to the
  $\Lp^\infty $ class.
\end{proof}

\begin{lemma}[Regularity and incompressibility of the velocity field on $\T^3$]
  \label{lem:global-velocity}
  \lean{FastDynamo.ScaleFamily.TorusCellPlacement.FullVelocityConcrete.parameterizedFullGlobalW1Inf,
    FastDynamo.ScaleFamily.TorusCellPlacement.FullVelocityConcrete.parameterizedFullGlobalW1Inf_isTransverseDivergenceFree,
    FastDynamo.ScaleFamily.TorusCellPlacement.FullVelocityConcrete.parameterizedFullGlobalW1Inf_ne_zero,
    FastDynamo.ScaleFamily.TorusCellPlacement.FullVelocityConcrete.lipschitzWith_literalVelocity_parameterizedFullGlobalRealTorusSource}
  \leanok
  \uses{def:velocity:global-velocity, lem:velocity:divergence-free, lem:velocity:w1inf-bound,
    lem:velocity:accumulation-decay, lem:velocity:nonvanishing}
  The field $u $ is a non-zero real-valued autonomous vector field belonging to
  $\W^{1,\infty}(\T^3;\R^3) $.  It is independent of $x_3 $, and $\Div u=0 $.
\end{lemma}

The paper's lemma ends with one further clause: \emph{$u $ is smooth away from
$\{p_\infty\}\times\T $ and is differentiable at every point of this circle, with derivative
zero.}  That clause is not part of the statement above, and it is not proved on $\T^3 $:
it is formalised for the chart-side representative and its periodisation in
Lemma~\ref{lem:main:velocity-regularity}.

Four remarks on the Lean form.  (a) Membership in $\W^{1,\infty} $ is carried by the type
\texttt{TransverseRealW1Inf}, so it is true by construction rather than proved; the
content is the construction itself --- that the infinite sum does define such an element
--- together with the explicit bound of Lemma~\ref{lem:velocity:w1inf-bound}.  The
$\W^{1,\infty} $ class lives on $\T^2 $ and reaches $\T^3 $ through the axially
independent extension \texttt{literalVelocity}, which is proved Lipschitz; there is no
separate $\W^{1,\infty}(\T^3) $ space in the development.  (b) Reality,
$x_3 $-independence and $\Div u=0 $ are Lemma~\ref{lem:velocity:divergence-free}.
(c) Non-vanishing holds for $\tau>0 $ (Lemma~\ref{lem:velocity:nonvanishing}).
(d) The smoothness and differentiability clauses are proved for the chart-side
representative and its $\Z^2 $-periodisation on the plane, not as smoothness in a manifold
structure on $\T^2 $ (the development installs none); they are collected, with the failure
of $\C^1 $, in the node \texttt{lem:main:velocity-regularity} and its Lean predicate
\texttt{GlobalVelocityQualitative}.

\begin{proof}
  \leanok
  \uses{lem:velocity:divergence-free, lem:velocity:w1inf-bound,
    lem:velocity:accumulation-decay, lem:velocity:nonvanishing,
    lem:velocity:disjoint-placement}
  Incompressibility and reality are Lemma~\ref{lem:velocity:divergence-free}, the uniform
  gradient bound is Lemma~\ref{lem:velocity:w1inf-bound}, non-vanishing is
  Lemma~\ref{lem:velocity:nonvanishing}, and differentiability with vanishing derivative on
  the accumulation circle is Lemma~\ref{lem:velocity:accumulation-decay}.  Away from
  $p_\infty $ the family of supports is locally finite, so the sum is locally a finite sum
  of smooth functions; with the previous sentence this makes $u $ differentiable at every
  point of $\T^3 $, with derivative bounded by $\norm{\nabla v}_\infty+\eta_\mathrm{ax} $
  (in Lean by the constant $L $ of Lemma~\ref{lem:velocity:w1inf-bound}).  Applying the mean value theorem to the periodic
  lift of $u $, we deduce that $u $ is globally Lipschitz with this bound.  Its classical
  divergence vanishes everywhere.  Since the weak derivative of a Lipschitz function agrees
  almost everywhere with its classical derivative, $u\in\W^{1,\infty} $ and $\Div u=0 $ in
  distributions.
\end{proof}

\textbf{Deviation: a Fourier dominated limit, not Rademacher.}  The route just narrated is
the paper's, and the two steps it rests on --- global
Lipschitz-ness from the mean value theorem on the periodic lift, and the vanishing of the
classical divergence --- are proved in Lean exactly that way
(\texttt{lipschitzWith\_parameterizedFullGlobalEuclideanRealSource}, obtained from
\texttt{norm\_\allowbreak fderiv\_\allowbreak parameterizedFullGlobalEuclideanRealSource\_\allowbreak le} through Mathlib's mean value inequality, and
\texttt{parameterizedFullGlobalTorusFirstDerivative\_transverseDivergence}, which is a
pointwise identity on the displayed representatives).  The last step is not: membership in
$\W^{1,\infty} $ is, in this development, membership in the closed Fourier graph
\texttt{TransverseRealW1Inf}, and no bridge \emph{Lipschitz $\Rightarrow$ Fourier graph}
--- a Rademacher theorem on $\T^2 $ together with an integration by parts for Lipschitz maps
--- exists in this library or in Mathlib in the required form.  Lean therefore builds the
graph element directly, by a route that needs neither a $\W^{1,1} $ summation of the cells
nor the paper's a.e.\ identification of weak and classical derivatives: the helical part is a dominated
limit of the Fourier coefficients of the partial sums (its cell derivatives are uniformly
bounded but \emph{not} $\Lp^\infty $-summable, so the graph norms cannot simply be added),
the axial part is an absolutely convergent $\Lp^\infty $ series, and the two graph elements
are added
(\texttt{integral\_parameterizedHelicalGlobal\_weakDerivative},
\texttt{parameterizedHelicalGlobalW1Inf}, \texttt{axialGlobalW1Inf} in
\texttt{FastDynamo/Global/TorusFullVelocityConcrete.lean}).  The conclusion --- that the
stored derivative fields are the Fourier weak derivatives of the stored value, and that
their diagonal sum vanishes --- is the paper's $u\in\W^{1,\infty} $, $\Div u=0 $.

\section{Why the construction stops at Lipschitz regularity}

The local field $v $ is smooth; regularity is lost only because the global field has
nonvanishing gradients at infinitely many shrinking horizontal scales.  The estimate of
Lemma~\ref{lem:velocity:accumulation-decay} shows that $u $ is nevertheless differentiable on
the accumulation circle, with $Du(p_\infty,x_3)=0 $ for every $x_3\in\T $.  At the other end,
both cut-offs are constant near each cell centre: the support fit puts $p_n $ well inside the
axial plateau $D(p_n,3R_n/4) $, and the rescaled local copy has amplitude $\ell_n $ and
dilation $\ell_n $, which cancel.  Hence, by Lemma~\ref{lem:velocity:matching-cell},
\[
  Du(p_n,x_3)=
  \begin{pmatrix}
    0      & -\Omega & 0 \\
    \Omega & 0       & 0 \\
    0      & 0       & 0
  \end{pmatrix}
  \neq0
  \qquad(x_3\in\T),
\]
that is, $Du(p_n,x_3)e_1=\Omega e_2 $, $Du(p_n,x_3)e_2=-\Omega e_1 $ and $Du(p_n,x_3)e_3=0 $;
it is non-zero because $\Omega\neq0 $.  It is one and the same map for every $n $, namely the
derivative at the origin of the unrescaled profile $v $.  Since $p_n\to p_\infty $, the
classical derivative $Du $ is discontinuous at every point of $\{p_\infty\}\times\T $.  A
continuous representative of the weak gradient would have to agree everywhere with $Du $, so
no such representative exists and
$u\notin\C^1 $, even though $u $ is differentiable everywhere (Section~3.4 of the paper).

These three statements --- differentiable everywhere, smooth off the accumulation point,
derivative discontinuous there --- are the content of the node
\texttt{lem:main:velocity-regularity} and of the Lean predicate
\texttt{GlobalVelocityQualitative} in
\texttt{FastDynamo/Global/TorusGlobalVelocityRegularity.lean}, and the Lean proof is the
argument just given, along the cell centres:
\texttt{fderiv\_\allowbreak parameterizedFullGlobalEuclideanRealSource\_\allowbreak cellCenter\_\allowbreak eq} gives one
and the same map at every centre,
\texttt{lineDeriv\_parameterizedSmoothedVelocity\_zero\_apply\_one} identifies its $(2,1) $
entry as the angular speed $\Omega $, whence
\texttt{fderiv\_\allowbreak parameterizedFullGlobalEuclideanRealSource\_\allowbreak cellCenter\_\allowbreak ne\_\allowbreak zero} under
$\Omega\neq0 $, and
\texttt{not\_\allowbreak continuousAt\_\allowbreak fderiv\_\allowbreak parameterizedFullGlobalEuclideanRealSource} concludes from
\texttt{tendsto\_cellCenter\_accumulationPoint} and
\texttt{fderiv\_\allowbreak parameterizedFullGlobalEuclideanRealSource\_\allowbreak accumulationPoint}, which says that
the derivative is $0 $ there.  Lean extracts only the entry $\Omega $ of the
matrix rather than the matrix itself, which is all the non-vanishing needs.  The global
failure of $\C^1 $ is the field \texttt{periodized\_not\_continuousAt\_fderiv} of the
strengthened main theorem (Theorem~\ref{thm:main-strengthened}): the derivative of the
$\Z^2$-periodised plane representative \texttt{mainVelocityPeriodizedRepresentative} is
discontinuous at the accumulation point
(\texttt{not\_\allowbreak continuousAt\_\allowbreak fderiv\_\allowbreak mainVelocityPeriodizedRepresentative}), a statement on
the plane, not on the manifold $\T^3 $ (Lemma~\ref{lem:main:velocity-regularity}).
